\documentclass{article}

\usepackage{amssymb,amsmath,amsthm,mathtools,cancel}
\usepackage[margin=1in]{geometry}
\usepackage{enumitem}


\newtheorem{theorem}{Theorem}[section]
\newtheorem{lemma}[theorem]{Lemma}

\begin{document}
\noindent
\begin{minipage}[t]{0.68\textwidth}
{\large 6.1200 \textit{Mathematics for Computer Science}}\\[4pt]
{\large Problem Set 2}\\[4pt]
{\large Solutions by Harsh Bhalala}
\end{minipage}%
\hfill
\begin{minipage}[t]{0.27\textwidth}
\raggedleft
{\large 18, September 2026}
\end{minipage}

\vspace{6pt}

\noindent\rule{\textwidth}{0.5pt}

\section*{Problem 1. Failed Induction}
\textbf{(a)}
    By taking n = 9, we get $3(9) - 6 = 21$ edges so we have 9 nodes and 21 edges.
    Now, in our counterexample, two of the 9 nodes have degree 0, and the remaining 7 nodes with 21 edges make
    a complete $K_7$ subgraph, which has $\chi(K_7) = 7$, thus proving the claim is false.

\vspace{\baselineskip}
\noindent\textbf{(b)} Not solved due to lack of Recitation.

\section*{Problem 2. High Chromatic Number Without Triangles}

\textbf{(a)}
    for G set of colors $c = $ \{red, blue, green, black\} and valid coloring function
    \[
    \begin{aligned}
    f = \{ &(w,\text{black}), (v_0,\text{blue}), (v_1,\text{green}), (v_2,\text{blue}), (v_3,\text{green}), (v_4,\text{red}),\\
    &(u_0,\text{blue}), (u_1,\text{red}), (u_2,\text{blue}), (u_3,\text{green}), (u_4,\text{red})
    \}
    \end{aligned}
    \]
Therefore, $\chi(G) \leq 4$.

\vspace{\baselineskip}
\noindent\textbf{(b)} \textbf{Lemma:} For a simple cicular graph $C$ with odd number of nodes, $\chi(C) = 3$.
\begin{proof}
    We know that a graph can only be two colored if it is a bipartite also a bipartie does not contain any
    odd-numbered cycle, meaning $C$ is not bipartite and cannot be 2-colored, but we can color nodes alternatively with
    2 colors and the last node in an odd position with one extra color, thus proving $\chi(C) = 3$.
\end{proof}

For graph $G$, the subgraph containing the outer nodes $v = {u_0,\dots,u_4}$ is odd numberd cicular graph, and from the lemma
we need at least 3 colors. Therefore, any valid coloring of $G$ must match the partial coloring shown.



\vspace{\baselineskip}
\noindent\textbf{(c)} By using part (b) and trying to color inner vertices $v'' = {v_0,\dots v_4}$, without loss of generality
    if we choose any three vertices from v'', suppose $v_1,v_4, v_3$, we observe that to color all three of them we at least need 3 colors
    and now to color the centeral vertex $w$ we must use 4th color, meaning $\chi(G) > 3$.
    three vertices

\section*{Problem 3. Find A Pet}

\textbf{(a)}
\begin{table}[h!]
\centering
\renewcommand{\arraystretch}{1.5}
\begin{tabular}{|c||l|l|l|l|l|l|l|}
\hline
& Day 1 & Day 2 & Day 3 & Day 4 & Day 5 & Day 6 & Day 7 \\
\hline
\hline
H & A, B, C, \fbox{D} & \fbox{D} & \fbox{D} & \fbox{D} & & & \\
\hline
I & $\emptyset$ & B, \fbox{C} & \fbox{C} & \fbox{C} & & & \\
\hline
S & $\emptyset$ & \fbox{A} & A, \fbox{B} & \fbox{B} & & & \\
\hline
Z & $\emptyset$ & $\emptyset$ & $\emptyset$ & \fbox{A} & & & \\
\hline
\end{tabular}
\end{table}

\pagebreak
\vspace{\baselineskip}
\noindent\textbf{(b)}

\begin{table}[h!]
\centering
\renewcommand{\arraystretch}{1.5}
\begin{tabular}{|c||l|l|l|l|l|l|l|}
\hline
& Day 1 & Day 2 & Day 3 & Day 4 & Day 5 & Day 6 & Day 7 \\
\hline
\hline
A & \fbox{I} & \fbox{I} & & & & & \\
\hline
B & $\emptyset$ & \fbox{Z} & & & & & \\
\hline
C & \fbox{S} & \fbox{S} & & & & & \\
\hline
D & \fbox{H}, Z & \fbox{H} & & & & & \\
\hline
\end{tabular}
\end{table}


\vspace{\baselineskip}
\noindent\textbf{(c)}
    Since H prefers D over C and C prefers H the best, they will form a rough matching.

\section*{Problem 4. Stable Merriment}

\textbf{(a)}
\begin{proof}
    Let's say for some arbitrary set of stable matchings between $n$ Pokemon and $n$ Trainers
    on the final day of the Stable Matching Algorithm, there are no merry Trainers. So at most each of the Trainers
    reject $n - \lceil \frac{n}{2}\rceil - 1$ Pokemon
    and since there are n Trainers
    \[S \leq n\left( n - \lceil \frac{n}{2}\rceil - 1 \right) < \frac{n^2}{2} \tag*{By elementary algebra}\]
\end{proof}

\vspace{\baselineskip}
\noindent\textbf{(b)}
    Let's say there are no merry Pokemon or Trainers. Since there are no merry Pokemon
    $S \geq n \lceil \frac{n}{2} \rceil \geq \frac{n^2}{2}$ but that would create a contradiction
    with part (a) meaning if there are no merry trainers then there must be at least one merry Pokemon
    and the opposite can be proved similarly

    Therefore, we conclude that at least one Pokemon or Trainer must be merry.
\end{document}